%
\vspace{-0.1cm}
\section{Introduction}
\label{sec:intro}

Deep neural networks provide rich representations that can enable reinforcement learning (RL) algorithms to perform effectively.
However, it was previously thought that the combination of simple online RL algorithms with deep neural networks was fundamentally unstable.
Instead, a variety of solutions have been proposed to stabilize the algorithm \citep{riedmiller2005nfq,mnih-atari-2013,mnih-dqn-2015,hado2015doubledqn,schulman2015trust}.
These approaches share a common idea: the sequence of observed data encountered by an online RL agent is non-stationary, and online RL updates are strongly correlated.
By storing the agent's data in an experience replay memory, the data can be batched \citep{riedmiller2005nfq,schulman2015trust} or randomly sampled~\citep{mnih-atari-2013,mnih-dqn-2015,hado2015doubledqn} from different time-steps.
Aggregating over memory in this way reduces non-stationarity and decorrelates updates, but at the same time limits the methods to off-policy reinforcement learning algorithms.

Deep RL algorithms based on experience replay have achieved unprecedented success in challenging domains such as Atari 2600.
However, experience replay has several drawbacks: it uses more memory and computation per real interaction; and it requires off-policy learning algorithms that can update from data generated by an older policy.

In this paper we provide a very different paradigm for deep reinforcement learning.
Instead of experience replay, we asynchronously execute multiple agents in parallel, on multiple instances of the environment.
This parallelism also decorrelates the agents' data into a more stationary process, since at any given time-step the parallel agents will be experiencing a variety of different states.
This simple idea enables a much larger spectrum of fundamental on-policy RL algorithms, such as Sarsa, n-step methods, and actor-critic methods, as well as off-policy RL algorithms such as Q-learning, to be applied robustly and effectively using deep neural networks.

Our parallel reinforcement learning paradigm also offers practical benefits.
%
Whereas previous approaches to deep reinforcement learning rely heavily on specialized hardware such as GPUs \citep{mnih-dqn-2015,hado2015doubledqn,schaul2015prioritized} or massively distributed architectures \citep{nair2015gorila}, our experiments run on a single machine with a standard multi-core CPU.
When applied to a variety of Atari 2600 domains, on many games asynchronous reinforcement learning achieves better results, in far less time than previous GPU-based algorithms, using far less resource than massively distributed approaches.
The best of the proposed methods, asynchronous advantage actor-critic (A3C), also mastered a variety of continuous motor control tasks as well as learned general strategies for exploring 3D mazes purely from visual inputs.
We believe that the success of A3C on both 2D and 3D games, discrete and continuous action spaces, as well as its ability to train feedforward and recurrent agents makes it the most general and successful reinforcement learning agent to date.
